


We consider the Uhlmann transformation algorithm in the purified query access model.
The formal statement of \Cref{infthm:uhlmenn purif query} is as follows, where $\sD$ is a $\big(\log{(d_\sA d_\sB)} + 1\big)$-qubit system. 
A quantum algorithm $\mathtt{UhlmannPurifiedQuery}$ is provided in \Cref{alg:Uhl purif query}, in which the symbol ``$\gets$" denotes assignment.
Recall that $s_{\rm min}$ and $r$ are the minimum non-zero singular value and the rank of $\sqrt{\sigma^\sA}\sqrt{\rho^\sA}$, respectively, and $U_\rho^{\hat{\sA}\hat{\sB}}$ and $U_\sigma^{\hat{\sA}\hat{\sB}}$ are unitaries that prepare states $\ket{\rho}^{\hat{\sA}\hat{\sB}}$ and $\ket{\sigma}^{\hat{\sA}\hat{\sB}}$, respectively.






\begin{theorem}[Uhlmann transformation algorithm in the purified query access model]
\label{thm:Uhlmann alg purif query model}
Let $\delta \in (0, 1)$ and $\chi \in \{\diamond, \tF\}$. Then, the quantum query algorithm $\mathtt{UhlmannPurifiedQuery}$ given by \Cref{alg:Uhl purif query} satisfies the following. 

A unitary $\til{W}_\diamond^{\sB\sD}$ given by $\til{W}_\diamond^{\sB\sD} = \mathtt{UhlmannPurifiedQuery}(U_\rho^{\hat{\sA}\hat{\sB}}, U_\sigma^{\hat{\sA}\hat{\sB}}; \delta, \diamond)$, satisfies that
\begin{equation}
\label{eq:error purif access}
    \f{1}{2}\big\|\til{\cW}_\diamond^{\sB\sD} \circ\cP_{\ket{0}}^{\bC\rarr\sD} - \til{\cU}_{\rm ideal}^{\sB\sD}\circ\cP_{\ket{0}}^{\bC\rarr\sD}\big\|_\diamond
    \leq \delta,
\end{equation}
where $\til{U}_{\rm ideal}^{\sB\sD}$ is an exact block-encoding unitary of the Uhlmann partial isometry $V^\sB$, and $\cP_{\ket{0}}^{\bC\rarr\sD}$ is the channel preparing the state $\ketbra{0}{0}^\sD$ on system $\sD$.
The algorithm uses $u_\diamond = \cO\big(\f{1}{s_{\rm min}}\log{\big(\f{1}{\delta}\big)}\big)$ queries to $U_\rho^{\hat{\sA}\hat{\sB}}$, $U_\sigma^{\hat{\sA}\hat{\sB}}$, and their inverses, where $s_{\min}$ is the minimum non-zero singular value of $\sqrt{\sigma^{\sA}}\sqrt{\rho^{\sA}}$ as in \Cref{tab:technical notation}.

Moreover, a unitary $\til{W}_{\tF}^{\sB\sD}$ given by $\til{W}_{\tF}^{\sB\sD} = \mathtt{UhlmannPurifiedQuery}(U_\rho^{\hat{\sA}\hat{\sB}}, U_\sigma^{\hat{\sA}\hat{\sB}}; \delta, \tF)$, satisfies that 
\begin{align}
        \rF\big(\cT^\sB(\ketbra{\rho}{\rho}^{\sA\sB}), \ket{\sigma}^{\sA\sB}\big) \geq \rF(\rho^\sA, \sigma^\sA) - \delta,
\end{align}
where $\cT^\sB = \tr_{\sD}\circ\til{\cW}_{\tF}^{\sB\sD}\circ\cP_{\ket{0}}^{\bC\rarr\sD}$.
The algorithm uses $u_{\tF} = \cO\Big(\min\big\{\f{1}{s_{\rm min}}, \f{r}{\delta}\big\}\log{\big(\f{1}{\delta}\big)}\Big)$ queries to $U_\rho^{\hat{\sA}\hat{\sB}}$, $U_\sigma^{\hat{\sA}\hat{\sB}}$, and their inverses, where $s_{\min}$ and $r$ are the minimum non-zero singular value and the rank of $\sqrt{\sigma^\sA}\sqrt{\rho^\sA}$ as in \Cref{tab:technical notation}.

In both cases, the quantum circuit for implementing the algorithm consists of $\cO\big(u_\chi \log{(d_\sA d_\sB)}\big)$ one- and two-qubit gates, and $\cO\big(\log{(d_\sA d_\sB)}\big)$ qubits suffice at any one time.
\end{theorem}




%===================================================================================================================================================================================================================================

\begin{algorithm}[h]
\caption{Uhlmann transformation algorithm in the purified query access model \\ \parbox{\linewidth}{\centering $\mathtt{UhlmannPurifiedQuery}(U_\rho^{\hat{\sA}\hat{\sB}}, U_\sigma^{\hat{\sA}\hat{\sB}}; \delta, \chi)$ (In Theorem~\ref{thm:Uhlmann alg purif query model})}}
\label{alg:Uhl purif query}
\SetKwInput{KwInput}{Input}
\SetKwInput{KwOutput}{Output}
\SetKwInput{KwParameters}{Parameters}

\SetAlgoNoEnd
\SetAlgoNoLine
\KwInput{Unitary oracles $U_\rho^{\hat{\sA}\hat{\sB}}$, $U_\sigma^{\hat{\sA}\hat{\sB}}$, and their inverses.}
\KwParameters{$\delta \in (0, 1)$ and $\chi \in \{\diamond, \tF\}$.}
\KwOutput{Unitary $\til{W}^{\sB\sD}$.}
\SetAlgoLined

Set $W^{\sB\hat{\sA}\hat{\sB}} \gets (U_\rho^{\hat{\sA}\hat{\sB}})^\dag(\bI^{\hat{\sA}} \otimes F^{\sB\hat{\sB}})U_\sigma^{\hat{\sA}\hat{\sB}}$, where $F^{\sB\hat{\sB}}$ is the swap operator between $\sB$ and $\hat{\sB}$. \\
\If{$\chi = \diamond$}{
  Set $\delta_1 \gets (\delta/3)^2$ and $\beta \gets s_{\rm min}$. \\
}
\ElseIf{$\chi = \tF$}{
  Set $\delta_1 \gets \delta/4$ and $\beta \gets \max\{s_{\rm min}, \delta_1/(2r)\}$. \\
}
Set $\til{W}^{\sB\sD} \gets \mathtt{QSVTSIGN}(W^{\sB\hat{\sA}\hat{\sB}}, (W^{\sB\hat{\sA}\hat{\sB}})^\dag; \delta_1, \beta)$ (Proposition~\ref{prop:FPAA general}).  \\
Return $\til{W}^{\sB\sD}$.
\end{algorithm}




As we show in the following subsection, the unitary
$W^{\sB\hat{\sA}\hat{\sB}}$ in Algorithm~\ref{alg:Uhl purif query} is a block-encoding of the operator
$M^{\sB} = \tr_{\sA'}\big[\ketbra{\sigma}{\rho}^{\sA'\sB}\big]$.
Applying $\mathtt{QSVTSIGN}$, the QSVT with the sign function described in \Cref{prop:FPAA general}, with input $W^{\sB\hat{\sA}\hat{\sB}}$ and $(W^{\sB\hat{\sA}\hat{\sB}})^\dag$ yields a unitary $\til{W}^{\sB\sD}$ that is a block-encoding of $P_\sgn^{(\rm SV)}(M^\sB)$. If the polynomial $P_\sgn$ sufficiently approximates the sign function, then by \Cref{prop:explicit form of Uhl}, $\til{W}^{\sB\sD}$ is close to a block-encoding unitary of the Uhlmann partial isometry $V^\sB$.



%===================================================================================================================================================================================================================================


We emphasize the importance of clarifying which measure is used to evaluate the degree of approximation to the Uhlmann transformation.
In Theorem~\ref{thm:Uhlmann alg purif query model}, the first statement evaluates the difference between the quantum \emph{channels} implemented by our algorithm and the ideal Uhlmann transformation using the diamond norm, accommodating the setting where input states to the Uhlmann transformation can be arbitrary.
On the other hand, the second statement evaluates the accuracy in terms of the output \emph{state} when the input to the Uhlmann transformation is $\ket{\rho}^{\sA\sB}$---where we measure the distance between the ideal and actual outputs by the fidelity.
While evaluating the approximation error via fidelity difference may suffice, it is indeed more convenient to use the diamond norm in cases where the Uhlmann transformation algorithm is used as a subroutine (e.g., in \Cref{sec:sample petz}).



One might think that post-selection is required to extract the Uhlmann partial isometry $V^\sB$ encoded in the top-left block of $U_{\rm ideal}^{\sB\sD}$. 
However, this is not the case, because the domain of the Uhlmann partial isometry covers the support of the state $\rho^\sB$, and thus, as long as the input state is $\ket{\rho}^{\sA\sB}$, the bottom-left block of $U_{\rm ideal}^{\sB\sD}$ does not affect the state: $U_{\rm ideal}^{\sB\sD}\ket{\rho}^{\sA\sB}\ket{0}^\sD = V^\sB\ket{\rho}^{\sA\sB}\ket{0}^{\sD}$.
This follows from the fact that $U_{\rm ideal}^{\sB\sD}$ is a unitary and $V^\sB$ is a partial isometry, which implies the bottom-left block acts non-trivially only on the kernel of $V^\sB$.





If we are able to implement $U_\rho^{\hat{\sA}\hat{\sB}}$, $U_\sigma^{\hat{\sA}\hat{\sB}}$, and their inverses by ourselves, the circuit complexity, a total number of one- and two-qubit gates, for the Uhlmann transformation can be quantified. 
This applies to both cases where the error is evaluated using the diamond norm and where it is evaluated using the fidelity difference.
Since $W^{\sB\hat{\sA}\hat{\sB}}$ is given by $(U_\rho^{\hat{\sA}\hat{\sB}})^\dag(\bI^{\hat{\sA}} \otimes F^{\sB\hat{\sB}})U_\sigma^{\hat{\sA}\hat{\sB}}$ (see line~1 of \Cref{alg:Uhl purif query}), its implementation requires the unitaries $(U_\rho^{\hat{\sA}\hat{\sB}})^\dag$ and $U_\sigma^{\hat{\sA}\hat{\sB}}$, along with the swap operation $F^{\sB\hat{\sB}}$, which can be realized using $\log d_\sB$ two-qubit swap gates.
Furthermore, the number of one- and two-qubit gates used in $\mathtt{QSVTSIGN}(W^{\sB\hat{\sA}\hat{\sB}}, (W^{\sB\hat{\sA}\hat{\sB}})^\dag; \delta_1, \beta)$ is $\cO\big(u_\chi\log{(d_\sA d_\sB)}\big)$, where $\chi \in \{\diamond, \tF\}$, because $W^{\sB\hat{\sA}\hat{\sB}}$ is a $(1, \log{(d_\sA d_\sB)}, 0)$-block-encoding unitary (see also Theorem~\ref{prop:general QSVT}). 
Therefore, the total number of one- and two-qubit gates for the overall algorithm is 
\begin{align}
    \cO\Big(u_\chi \big(\cC(U_\rho) + \cC(U_\sigma) + \log{d_\sB}\big) + u_\chi\log{(d_\sA d_\sB)} \Big) 
    =\cO\big(u_\chi\big(\cC(U_\rho) + \cC(U_\sigma) + \log{(d_\sA d_\sB)}\big)\big),
\end{align}
where $\cC(U_\rho)$ and $\cC(U_\sigma)$ are the number of one- and two-qubit gates required to implement $U_\rho^{\hat{\sA}\hat{\sB}}$ and $U_\sigma^{\hat{\sA}\hat{\sB}}$, respectively.





We first prove Theorem~\ref{thm:Uhlmann alg purif query model} in \Cref{sec:proof purif query Uhl}, and then provide a proof of a technical lemma used to show the theorem in \Cref{sec:proof of lem FPAA diamond}.






%=============================================================================================================================

\subsection{Proof of Theorem~\ref{thm:Uhlmann alg purif query model}}
\label{sec:proof purif query Uhl}



Let $U_\rho^{\hat{\sA}\hat{\sB}}$ and $U_\sigma^{\hat{\sA}\hat{\sB}}$ be unitaries such that $U_\rho^{\hat{\sA}\hat{\sB}}\ket{0}^{\hat{\sA}\hat{\sB}} = \ket{\rho}^{\hat{\sA}\hat{\sB}}$, and $U_\sigma^{\hat{\sA}\hat{\sB}}\ket{0}^{\hat{\sA}\hat{\sB}} = \ket{\sigma}^{\hat{\sA}\hat{\sB}}$.
We then define a unitary $W^{\sB\hat{\sA}\hat{\sB}}$ by $W^{\sB\hat{\sA}\hat{\sB}} = (U_\rho^{\hat{\sA}\hat{\sB}})^\dag(\bI^{\hat{\sA}} \otimes F^{\sB\hat{\sB}})U_\sigma^{\hat{\sA}\hat{\sB}}$.
The unitary $W^{\sB\hat{\sA}\hat{\sB}}$ is a $(1, \log{(d_\sA d_\sB)}, 0)$-block-encoding of $\tr_{\sA'}\big[\ketbra{\sigma}{\rho}^{\sA'\sB}\big]$:
\begin{align}
\label{eq:block W fidelity}
        W^{\sB\hat{\sA}\hat{\sB}}
        = \hspace{1mm}
\begin{blockarray}{ccc}
& \bra{0}^{\hat{\sA}\hat{\sB}}&  & \vspace{1mm}\\
\begin{block}{c(cc)}
  \ket{0}^{\hat{\sA}\hat{\sB}} \hspace*{1mm} & \tr_{\sA'}\big[\ketbra{\sigma}{\rho}^{\sA'\sB}\big] & \hspace{0mm} * \hspace{3mm}\\
   \hspace*{1mm} & * & \hspace{0mm} * \hspace{3mm}\\
\end{block}
\end{blockarray}\hspace{2.5mm}.
\end{align}
To see this, we use the Schmidt decomposition of $\ket{\rho}^{\sA\sB}$ and $\ket{\sigma}^{\sA\sB}$:
\begin{align}
      \ket{\rho}^{\sA\sB} = \sum_{k=1}^{r_\rho}\sqrt{p_k}\ket{e_k}^\sA\ket{f_k}^\sB, \ \ \text{and} \ \ \ 
      \ket{\sigma}^{\sA\sB} = \sum_{l=1}^{r_\sigma}\sqrt{q_l}\ket{g_l}^\sA\ket{h_l}^\sB.
\end{align}
We then compute the elements of $W^{\sB\hat{\sA}\hat{\sB}}$ as
\begin{align}
    \bra{i}^\sB\bra{0}^{\hat{\sA}\hat{\sB}}W^{\sB\hat{\sA}\hat{\sB}}\ket{0}^{\hat{\sA}\hat{\sB}}\ket{j}^{\sB}
    &=  \bra{i}^\sB\bra{\rho}^{\hat{\sA}\hat{\sB}}(\bI^{\hat{\sA}} \otimes F^{\sB\hat{\sB}})\ket{\sigma}^{\hat{\sA}\hat{\sB}}\ket{j}^{\sB} \\
    &= \sum_{k, l} \sqrt{p_k}\sqrt{q_l}\braket{e_k}{g_l}^{\hat{\sA}} \bra{i}^\sB\bra{f_k}^{\hat{\sB}} F^{\sB\hat{\sB}} \ket{h_l}^{\hat{\sB}}\ket{j}^\sB \\
    &= \sum_{k, l} \sqrt{p_k}\sqrt{q_l}\braket{e_k}{g_l}\braket{i}{h_l}\braket{f_k}{j} \\
    &= \bra{i}^\sB \Big(\sum_{k, l} \sqrt{p_k}\sqrt{q_l}\braket{e_k}{g_l}\ketbra{h_l}{f_k}^\sB \Big) \ket{j}^\sB \\
    \label{inteq:171}
    &= \bra{i}^\sB \big(\tr_{\sA'}\big[\ketbra{\sigma}{\rho}^{\sA'\sB}\big]\big) \ket{j}^\sB.
\end{align}
Thus, $\bra{0}^{\hat{\sA}\hat{\sB}}W^{\sB\hat{\sA}\hat{\sB}}\ket{0}^{\hat{\sA}\hat{\sB}} = \tr_{\sA'}\big[\ketbra{\sigma}{\rho}^{\sA'\sB}\big]$.



Due to \Cref{prop:explicit form of Uhl}, it suffices to apply the sign function on the top-left block of $W^{\sB\hat{\sA}\hat{\sB}}$.
We implement the QSVT with the sign function $\mathtt{QSVTSIGN}$ in \Cref{prop:FPAA general} to approximate the sign function by a polynomial $P_\sgn$.
We denote by $\til{W}^{\sB\sD}$ a unitary which is a block-encoding of $P_\sgn^{(\rm SV)}\big(\tr_{\sA'}\big[\ketbra{\sigma}{\rho}^{\sA'\sB}\big]\big)$:
\begin{align}
\label{eq:block tildeW fidelity}
    \til{W}^{\sB\sD}
    = \hspace{1mm}
\begin{blockarray}{ccc}
& \bra{0}^{\sD}&  & \vspace{1mm}\\
\begin{block}{c(cc)}
  \ket{0}^{\sD} \hspace*{1mm} & P_\sgn^{(\rm SV)}\big(\tr_{\sA'}\big[\ketbra{\sigma}{\rho}^{\sA'\sB}\big]\big) & \hspace{0mm} * \hspace{3mm}\\
   \hspace*{1mm} & * & \hspace{0mm} * \hspace{3mm}\\
\end{block}
\end{blockarray}\hspace{2.5mm}, 
\end{align}
where $\sD$ is a system including $\log{(d_\sA d_\sB)} + 1$ qubits.



The discussion so far applies to both cases, in which the performance of our algorithm is measured in terms of the diamond norm and the fidelity difference.
In the following, we discuss each case separately.
In Secs.~\ref{sec:eval diamond in purif query} and~\ref{sec:uhlfidquery}, we evaluate the approximation error using the diamond norm and the fidelity difference, respectively.
These results demonstrate the complete statement of Theorem~\ref{thm:Uhlmann alg purif query model}.


%================================================================================================================================



\subsubsection{Evaluation in the diamond norm}
\label{sec:eval diamond in purif query}



To evaluate the error in the diamond norm, we use the following lemma.
For its application in the following sections, we state it in a slightly more general form than the form used in this section. We provide a proof of this lemma in \Cref{sec:proof of lem FPAA diamond}.

\begin{lemma}
\label{lem:FPAA diamond}
    Let $\delta \in (0, 1)$, $M^\sB$ be a matrix, and $\ket{\varphi}^\sE$ and $\ket{\phi}^\sE$ be two pure states. 
    Suppose that $\til{U}^{\sB\sD\sE}$ is a block-encoding unitary of $\ketbra{\varphi}{\phi}^\sE\otimes P_\sgn^{(\rm SV)}(M^\sB)$ that satisfies
\begin{align}
\label{inteq:24}
    \big\|P_\sgn^{(\rm SV)}(M^\sB) - \sgn^{(\rm SV)}(M^\sB)\big\|_\infty \leq \delta.
\end{align}
Then, there exists a unitary $\til{U}_{\rm ideal}^{\sB\sD\sE}$ which is an exact block-encoding of $\ketbra{\varphi}{\phi}^\sE \otimes \sgn^{(\rm SV)}(M^\sB)$ and satisfies that 
\begin{align}
\label{eq:distance of L and L_ideal}
    \f{1}{2}\big\|\til{\cU}^{\sB\sD\sE} \circ \cP_{\ket{0}\ket{\phi}}^{\bC\rarr\sD\sE} - \til{\cU}_{\rm ideal}^{\sB\sD\sE} \circ \cP_{\ket{0}\ket{\phi}}^{\bC\rarr\sD\sE}\big\|_\diamond \leq 3\sqrt{\delta},
\end{align}
where $\cP_{\ket{0}\ket{\phi}}^{\bC\rarr\sD\sE}$ is a state-preparation channel such that $\cP_{\ket{0}\ket{\phi}}^{\bC\rarr\sD\sE} = \ketbra{0}{0}^{\sD}\otimes\ketbra{\phi}{\phi}^\sE$.
\end{lemma}





When we use $\mathtt{QSVTSIGN}$ in \Cref{prop:FPAA general}, we set $\beta$ to the minimum non-zero singular value of $\tr_{\sA'}\big[\ketbra{\sigma}{\rho}^{\sA'\sB}\big]$, denoted by $s_{\rm min}$.
Then, we obtain a block-encoding unitary 
\begin{align}
    \til{W}^{\sB\sD} = \mathtt{QSVTSIGN}(W^{\sB\hat{\sA}\hat{\sB}}, (W^{\sB\hat{\sA}\hat{\sB}})^\dag; \delta_1, \beta_\diamond)    
\end{align}
of a polynomial $P_\sgn^{(\rm SV)}\big(\tr_{\sA'}\big[\ketbra{\sigma}{\rho}^{\sA'\sB}\big]\big)$, which satisfies  
\begin{align}
\label{inteq:65}
    \big\|P_\sgn^{(\rm SV)}\big(\tr_{\sA'}\big[\ketbra{\sigma}{\rho}^{\sA'\sB}\big]\big) - \sgn^{(\rm SV)}\big(\tr_{\sA'}\big[\ketbra{\sigma}{\rho}^{\sA'\sB}\big]\big)\big\|_\infty \leq \delta_1,
\end{align}
where $\beta_\diamond = s_{\rm min}$ and $\delta_1 \in (0, 1/2)$.
The number of iteration of the unitary $W^{\sB\hat{\sA}\hat{\sB}}$ is given by the minimum odd integer $u_\diamond$ satisfying $u_\diamond \geq \big\lceil\f{8e}{\beta_\diamond}\log{(2/\delta_1)}\big\rceil = \big\lceil\f{8e}{s_{\rm min}}\log{(2/\delta_1)}\big\rceil$.
By direct calculation, one can check that $s_{\rm min}$ corresponds to the minimum non-zero singular value of $\sqrt{\sigma^\sA}\sqrt{\rho^\sA}$ (see also \Cref{sec:analyze partial iso}).

From Eq.~\eqref{inteq:65} and \Cref{lem:FPAA diamond} with $d_\sE = 1$, we have that
\begin{align}
    \f{1}{2}\big\|\til{\cW}^{\sB\sD}\circ\cP_{\ket{0}}^{\bC\rarr\sD} - \cU_{\rm ideal}^{\sB\sD}\circ \cP_{\ket{0}}^{\bC\rarr\sD}\big\|_\diamond \leq 3\sqrt{\delta_1},
\end{align}
where $U_{\rm ideal}^{\sB\sD}$ is an exact block encoding of $\sgn^{(\rm SV)}\big(\tr_{\sA'}\big[\ketbra{\sigma}{\rho}^{\sA'\sB}\big]\big)$, which is nothing but the Uhlmann partial isometry $V^\sB$ (see \Cref{prop:explicit form of Uhl}).
As the number of queries to $W^{\sB\hat{\sA}\hat{\sB}}$ and $(W^{\sB\hat{\sA}\hat{\sB}})^\dag$ are given by $u_\diamond = \cO\big(\log{(1/\delta_1)}/s_{\rm min}\big)$, the number of queries to $U_\rho^{\hat{\sA}\hat{\sB}}$, $(U_\rho^{\hat{\sA}\hat{\sB}})^\dag$, $U_\sigma^{\hat{\sA}\hat{\sB}}$, and $(U_\sigma^{\hat{\sA}\hat{\sB}})^\dag$ are of the same.
By rescaling $\delta_1$ as $\delta_1 = (\delta/3)^2$, we obtain the result for the number of queries.
Since this algorithm is conducted sequentially, $\cO\big(\log{(d_\sA d_\sB)}\big)$ qubits are needed at once.









%====================================================================================================================================





\subsubsection{Evaluation in the fidelity difference}
\label{sec:uhlfidquery}



When we evaluate the error between the ideal Uhlmann transformation and the transformation $\cT^\sB$ realized by our algorithm using the fidelity difference,
\begin{align}
    \rF(\rho^\sA, \sigma^\sA) - \rF(\cT^\sB(\ketbra{\rho}{\rho}^{\sA\sB}), \ket{\sigma}^{\sA\sB}),
\end{align}
there is a case where the transformation can be realized with fewer queries than $\cO\big(\log{(1/\delta)/s_{\rm min}}\big)$ derived in the previous section.


Let $\til{W}^{\sB\sD} = \mathtt{QSVTSIGN}(W^{\sB\hat{\sA}\hat{\sB}}, (W^{\sB\hat{\sA}\hat{\sB}})^\dag; \delta_1, \beta)$, and $M^\sB = \tr_{\sA'}\big[\ketbra{\sigma}{\rho}^{\sA'\sB}\big]$.
The square root fidelity between $\til{W}^{\sB\sD}\ket{\rho}^{\sA\sB}\ket{0}^{\sD}$ and $\ket{\sigma}^{\sA\sB}\ket{0}^{\sD}$ is given by 
\begin{align}
\label{inteq:14}
    \sqrt{\rF'} &=\sqrt{\rF}\big(\til{W}^{\sB\sD}\ket{\rho}^{\sA\sB}\ket{0}^{\sD}, \ket{\sigma}^{\sA\sB}\ket{0}^{\sD}\big) \\
    &= \big|\bra{\sigma}^{\sA\sB}\bra{0}^{\sD}\til{W}^{\sB\sD}\ket{0}^{\sD}\ket{\rho}^{\sA\sB}\big| \\
    &= \big|\bra{\sigma}^{\sA\sB}P_\sgn^{(\rm SV)}(M^\sB)\ket{\rho}^{\sA\sB}\big|,
\end{align}
where $P_\sgn$ is a polynomial such that $|P_\sgn(x) - \sgn(x)| \leq \delta_1$ for $|x| \in [\beta, 1]$ and $|P_\sgn(x)| \leq 1$ for $|x| \in [0, 1]$.
We then evaluate the difference between $\sqrt{\rF'}$ and $\sqrt{\rF} = \sqrt{\rF}(\rho^\sA, \sigma^\sA) = \sqrt{\rF}(V^\sB\ket{\rho}^{\sA\sB}, \ket{\sigma}^{\sA\sB})$ as
\begin{align}
    \label{inteq:104}
    \big|\sqrt{\rF} - \sqrt{\rF'}\big| 
    &= \big||\bra{\sigma}^{\sA\sB}V^\sB\ket{\rho}^{\sA\sB}| - |\bra{\sigma}^{\sA\sB}P_\sgn^{(\rm SV)}(M^\sB)\ket{\rho}^{\sA\sB}| \big| \\
    &\leq \big| \bra{\sigma}^{\sA\sB}V^\sB\ket{\rho}^{\sA\sB} - \bra{\sigma}^{\sA\sB}P_\sgn^{(\rm SV)}(M^\sB)\ket{\rho}^{\sA\sB} \big| \\
    &= \Big|\tr\big[\big(V^\sB - P_\sgn^{(\rm SV)}(M^\sB)\big)\tr_{\sA}\big[\ketbra{\rho}{\sigma}^{\sA\sB}\big]\big]\Big| \\
    &= \Big|\tr\big[\big(\sgn^{(\rm SV)}(M^\sB) - P_\sgn^{(\rm SV)}(M^\sB)\big)(M^{\sB})^\dag\big]\Big| \\
    &= \Big|\sum_k \big(1 - P_\sgn(s_k)\big)s_k\Big| \\
    \label{inteq:8}
    &\leq \sum_k \big|1 - P_\sgn(s_k)\big|s_k,
\end{align}
where $\{s_k\}_k$ are the singular values of $M^\sB$.
We should note that they correspond to the singular values of $\sqrt{\sigma^\sA}\sqrt{\rho^\sA}$ and that $\sum_k s_k = \big\|\sqrt{\sigma^\sA}\sqrt{\rho^\sA}\big\|_1 = \sqrt{\rF} \leq 1$.



We divide the label $k$ of the singular values into two parts $I_\beta$ and $\bar{I}_\beta$, where $I_\beta = \{k \in \bN; s_k \geq \beta\}$ and $\bar{I}_\beta = \{k \in \bN; s_k < \beta\}$.
Then, we compute Eq.~\eqref{inteq:8} as
\begin{align}
    \sum_k \big|1 - P_\sgn(s_k)\big|s_k
    &= \sum_{k\in I_\beta} \big|1 - P_\sgn(s_k)\big|s_k 
    + \sum_{k\in \bar{I}_\beta} \big|1 - P_\sgn(s_k)\big|s_k \\
    &\leq \delta_1 \sum_{k\in I_\beta} s_k + 2\sum_{k\in \bar{I}_\beta}s_k, \\
    &\leq \delta_1 + 2\beta \#\bar{I}_\beta \\
    &\leq \delta_1 + 2\beta r,
\end{align}
where $\#\bar{I}_\beta$ is the number of elements in $\bar{I}_\beta$ and $r$ is the rank of $\sqrt{\sigma^\sA}\sqrt{\rho^\sA}$.
We used \Cref{prop:FPAA general} and the inequality $\big|1 - P_\sgn(x)\big| \leq 2$ for $x \in [0, 1]$.
We also used the fact that $\#\bar{I}_\beta$ is at most $r$.


When we choose the threshold $\beta$, as $1/\beta = 2r/\delta_1$, we have that
$\big|\sqrt{\rF} - \sqrt{\rF'}| \leq 2\delta_1$.
On the other hand, if we set $1/\beta = 1/s_{\min}$, then $\#\bar{I}_\beta$ becomes zero, which leads to $\big|\sqrt{\rF} - \sqrt{\rF'}\big| \leq \delta_1$.
Thus, $\beta$ is chosen as $1/\beta =  1/\beta_\tF = \min\{1/s_{\rm min}, 2r/\delta_1\}$ to ensure at least 
\begin{align}
\label{inteq:13}
    \big|\sqrt{\rF} - \sqrt{\rF'}\big| \leq 2\delta_1.
\end{align}
Using the inequality that $|x-y|\leq 2 |\sqrt{x} - \sqrt{y}|$ for $0 \leq x, y \leq 1$, we have that
\begin{align}
    \big|\rF\big(\til{W}^{\sB\sD}\ket{\rho}^{\sA\sB}\ket{0}^{\sD}, \ket{\sigma}^{\sA\sB}\ket{0}^{\sD}\big)
    - \rF(\rho^\sA, \sigma^\sA)\big|\leq 4 \delta_1.
\end{align}


By denoting $\tr_{\sD}\circ\til{\cW}^{\sB\sD}\circ\cP_{\ket{0}}^{\bC\rarr\sD}$ by $\cT^\sB$ and rescaling $\delta_1$ as $\delta_1 = \delta/4$, we obtain that
\begin{align}
    \rF(\rho^\sA, \sigma^\sA) - \delta 
    &\leq \rF\big(\til{W}^{\sB\sD}\ket{\rho}^{\sA\sB}\ket{0}^{\sD}, \ket{\sigma}^{\sA\sB}\ket{0}^{\sD}\big) \\
    &\leq \rF\big(\cT^\sB(\ketbra{\rho}{\rho}^{\sA\sB}), \ket{\sigma}^{\sA\sB}\big),
\end{align}
where we used the monotonicity of the fidelity under the partial trace. 
The number of queries to $U_\rho^{\hat{\sA}\sB}$, $U_\sigma^{\hat{\sA}\sB}$, and their inverses is evaluated as
\begin{align}
    u_\tF &= \cO\big(\log{(1/\delta)}/\beta_\tF\big) \\
    &= \cO\Big(\min\Big\{\f{1}{s_{\rm min}}, \f{r}{\delta}\Big\}\log{\Big(\f{1}{\delta}\Big)}\Big),
\end{align}
and at any one time, $\cO\big(\log(d_\sA d_\sB)\big)$ ancilla qubits suffice.
We complete a proof of Theorem~\ref{thm:Uhlmann alg purif query model}.











%=================================================================================================================================



\subsection{Proof of \Cref{lem:FPAA diamond}}
\label{sec:proof of lem FPAA diamond}




We now give the proof of \Cref{lem:FPAA diamond}.
\begin{proof}[Proof of \Cref{lem:FPAA diamond}]

We use the (right) polar decomposition; for any matrix $A$, there exists a partial isometry $V$ with $\supp[V] \supseteq \im[A^\dag A]$ such that $A$ can be expressed as $A = V\sqrt{A^\dag A}$.
Since $\til{U}^{\sB\sD\sE}$ is a unitary, we have $(\til{U}^{\sB\sD\sE})^\dag \til{U}^{\sB\sD\sE} = \bI^{\sB\sD\sE}$.
Thus, by applying the polar decomposition to the bottom-left block of $\til{U}^{\sB\sD\sE}$, its first column block can be represented as
\begin{align}
\til{U}^{\sB\sD\sE} =
\begin{blockarray}{ccc}
& \bra{0}^{\sD}&  \vspace{1mm}\\
\begin{block}{c(cc)} 
  \ket{0}^{\sD}\hspace*{1mm} &  \ketbra{\varphi}{\phi}^\sE\otimes P_\sgn^{(\rm SV)}(M^\sB) & \hspace*{2mm} * \hspace*{2mm}\\ 
  &  V^{\sB\sE \rarr \sB\sD\sE}\sqrt{\bI^{\sB\sE} - S^{\sB\sE}} &\hspace*{2mm} * \hspace*{2mm}\\
\end{block}
\end{blockarray}\hspace{2.5mm},
\end{align}
where $S^{\sB\sE} = \ketbra{\phi}{\phi}^\sE \otimes P_\sgn^{(\rm SV)}((M^\sB)^\dag)P_\sgn^{(\rm SV)}(M^\sB)$, and $V^{\sB\sE \rarr \sB\sD\sE}$ is a partial isometry such that $\supp\big[V^{\sB\sE\rarr\sB\sD\sE}\big] \supseteq \im\big[\bI^{\sB\sE} - S^{\sB\sE}\big]$.
We define a unitary $\til{U}_{\rm ideal}^{\sB\sD\sE}$, which is exactly block-encoding of $\ketbra{\varphi}{\phi}^\sE \otimes \sgn^{(\rm SV)}(M^\sB)$, by
\begin{align}
\til{U}_{\rm ideal}^{\sB\sD\sE} = 
\begin{blockarray}{ccc}
& \bra{0}^{\sD}&  \vspace{1mm}\\
\begin{block}{c(cc)} 
  \ket{0}^{\sD}\hspace*{1mm} &  \ketbra{\varphi}{\phi}^\sE\otimes\sgn^{(\rm SV)}(M^\sB) & \hspace*{2mm} * \hspace*{2mm}\\
  &  V^{\sB\sE \rarr \sB\sD\sE}\Pi_\perp^{\sB\sE} &\hspace*{2mm} * \hspace*{2mm}\\
\end{block}
\end{blockarray}\hspace{2.5mm},
\end{align}
where $\Pi_\perp^{\sB\sE}$ is a projection onto $\ker[\ketbra{\phi}{\phi}^\sE\otimes M^\sB]$, given by 
\begin{align}
    \Pi_\perp^{\sB\sE} = \bI^{\sB\sE} - \ketbra{\phi}{\phi}^\sE\otimes \sgn^{(\rm SV)}((M^\sB)^\dag) \sgn^{(\rm SV)}(M^\sB).
\end{align}

For any state $\ket{\psi}^{\sR\sB}$ with any size of reference system $\sR$, we consider two states such that
\begin{align}
    \ket{\Psi}^{\sR\sB\sD\sE} 
    &= \til{U}^{\sB\sD\sE}\ket{0}^{\sD}\ket{\phi}^\sE\ket{\psi}^{\sR\sB} \\
    &=\ket{0}^{\sD}\ket{\phi}^\sE\otimes P_\sgn^{(\rm SV)}(M^\sB)\ket{\psi}^{\sR\sB} 
    + V^{\sB\sE \rarr \sB\sD\sE}\sqrt{\bI^{\sB\sE} - S^{\sB\sE}}\ket{\phi}^\sE\ket{\psi}^{\sR\sB},
\end{align}
and
\begin{align}
    \ket{\Psi'}^{\sR\sB\sD\sE} 
    &= \til{U}_{\rm ideal}^{\sB\sD\sE}\ket{0}^{\sD}\ket{\phi}^\sE\ket{\psi}^{\sR\sB} \\
    &= \ket{0}^{\sD}\ket{\phi}^\sE \otimes \sgn^{(\rm SV)}(M^\sB)\ket{\psi}^{\sR\sB} 
    + V^{\sB\sE \rarr \sB\sD\sE}\Pi_\perp^{\sB\sE}\ket{\phi}^\sE\ket{\psi}^{\sR\sB}.
\end{align}
The distance between these states is bounded as
\begin{align}
\label{inteq:5}
    \big\|\ket{\Psi}^{\sR\sB\sD\sE} - \ket{\Psi'}^{\sR\sB\sD\sE}\big\| 
    \leq \big\|P_\sgn^{(\rm SV)}(M^\sB) - \sgn^{(\rm SV)}(M^\sB)\big\|_\infty
    + \Big\|\sqrt{\bI^{\sB\sE} - S^{\sB\sE}} - \Pi_\perp^{\sB\sE}\Big\|_\infty.
\end{align}
The first term in the right-hand side of Eq.~\eqref{inteq:5} is, by assumption Eq.~\eqref{inteq:24}, bounded above by $\delta \in (0, 1)$. 
Regarding the second term, we can see that
\begin{align}
    \sqrt{\bI^{\sB\sE} - S^{\sB\sE}}
    &= \sqrt{\bI^{\sB\sE} - \ketbra{\phi}{\phi}^\sE\otimes P_\sgn^{(\rm SV)}((M^\sB)^\dag)P_\sgn^{(\rm SV)}(M^\sB)} \\
    &= \sqrt{\sum_{k=1}^r\Big(1-\big(P_\sgn(s_k)\big)^2\Big)\ketbra{\phi}{\phi}^\sE\otimes\ketbra{\xi_k}{\xi_k}^\sB + \Pi_\perp^{\sB\sE}} \\
    \label{inteq:4}
    &= \sum_{k=1}^r\sqrt{1-\big(P_\sgn(s_k)\big)^2}\ketbra{\phi}{\phi}^\sE\otimes\ketbra{\xi_k}{\xi_k}^{\sB} + \Pi_\perp^{\sB\sE},
\end{align}
where the singular value decomposition of $M^\sB$ is expressed by $\sum_{k=1}^{r}s_k\ketbra{\eta_k}{\xi_k}^\sB$.
Thus, the second term in the right-hand side of Eq.~\eqref{inteq:5} is bounded as 
\begin{align}
    \Big\|\sqrt{\bI^{\sB\sE} -S^{\sB\sE}} - \Pi_\perp^{\sB\sE}\Big\|_\infty 
    &=\max_{k \in [1, r]}\Big|\sqrt{1-\big(P_\sgn(s_k)\big)^2}\Big| \\
    &\leq\max_{k \in [1, r]}\Big|\sqrt{2\big(1-P_\sgn(s_k)\big)}\Big| \\
    &\leq \sqrt{2\delta}.
\end{align}
In the first inequality, we used $1-x^2 \leq 2(1-x)$ for any $x$, and in the last inequality, we used the assumption Eq.~\eqref{inteq:24}.


Hence, we can bound Eq.~\eqref{inteq:5} as
\begin{equation}
\label{inteq:20}
    \big\|\ket{\Psi}^{\sR\sB\sD\sE} - \ket{\Psi'}^{\sR\sB\sD\sE}\big\| \leq \delta + \sqrt{2\delta}.
\end{equation}
Using the fact that $\f{1}{2}\|\ketbra{\Psi}{\Psi}^{\sR\sB\sD\sE} - \ketbra{\Psi'}{\Psi'}^{\sR\sB\sD\sE}\|_1 \leq \|\ket{\Psi}^{\sR\sB\sD\sE} - \ket{\Psi'}^{\sR\sB\sD\sE}\|$ (see Eq.~\eqref{eq:relation of Euclidean and trace norm}) and that Eq.~\eqref{inteq:20} holds for any state $\ket{\psi}^{\sR\sB}$ with any system $\sR$, we can conclude that
\begin{align}
    \f{1}{2}\big\|\til{\cU}^{\sB\sD\sE} \circ \cP_{\ket{0}\ket{\phi}}^{\bC\rarr\sD\sE} - \til{\cU}_{\rm ideal}^{\sB\sD\sE} \circ \cP_{\ket{0}\ket{\phi}}^{\bC\rarr\sD\sE}\big\|_\diamond &\leq \delta + \sqrt{2\delta} \\
    \label{inteq:15}
    &\leq 3\sqrt{\delta},
\end{align}
where $\delta \in (0, 1)$ and $\cP_{\ket{0}\ket{\phi}}^{\bC\rarr\sD\sE} = \ketbra{0}{0}^\sD \otimes \ketbra{\phi}{\phi}^\sE$.

\end{proof}

